% ============================================================
%  THE ORTYX PROTOCOL
%  Part III — Decomposition
%  Chapter 13: The Four Levels Disentangled
% ============================================================

\chapter{The Four Levels Disentangled}
\label{ch:four-levels}

\chapepigraph{The most common error in speculative theoretical work
is not claiming too much but conflating levels. A claim at Level~II
dressed in Level~IV language is not a metaphysical discovery;
it is a useful analogy that has forgotten it is an analogy.}{Flyxion}

\noindent
The four-level taxonomy introduced in Chapter~2 and deployed
throughout Parts~I and~II provides the second major axis
of the decomposition. Where the failure geometry analysis
asks \textit{how} claims fail to provide epistemic grounding,
the four-level analysis asks \textit{what kind of claim}
is being made. The two axes are complementary: a claim
at \LevelIV{} exhibiting $\FailGeo{2}$ is different in
its remediation requirements from a claim at \LevelI{}
exhibiting $\FailGeo{1}$.

The four levels are:

\LevelI{}: Engineering utility. Claims about measurable
performance advantages of specific implementations.

\LevelII{}: Computational metaphor. Claims about the
structural analogy between the framework's mechanisms
and phenomena in other domains, where the analogy
illuminates without asserting identity.

\LevelIII{}: Cognitive phenomenology. Claims about
the character of experience, about what it is like
to perceive, remember, or attend in ways consistent
with the framework's account.

\LevelIV{}: Ontological claim. Claims about what
cognitive phenomena, conscious experience, or meaningful
structure fundamentally \textit{are} --- claims that
assert not just that the framework provides a useful
model but that the framework captures the deep structure
of the phenomena it addresses.

\section{Level~I Claims: The Stable Core}
\label{sec:level1-claims}

\LevelI{} claims in the Phoenix corpus are its most
stable. They include:

\textit{Jitter metrics:} The claim that Marine-derived
jitter is higher in block-coded compressed audio than
in the original waveform at block boundaries. This is
a specific, testable prediction about a measurable
quantity.

\textit{Compression efficiency:} The claim that harmonic
consensus encoding achieves lower bit rates than
independent harmonic encoding on strongly coupled
sources. This follows from the conformity relation and
is testable on any dataset of harmonic audio.

\textit{O(1) salience estimation:} The claim that the
Marine algorithm achieves salience estimates at constant
time per peak without global frequency analysis. This
is a computational complexity claim, verifiable by
implementation.

\textit{Associative memory efficiency:} The claim that
wave superposition encodes $O(N^2)$ pairwise associations
at $O(N)$ storage cost. This is a mathematical claim,
verifiable by derivation.

\textit{Context sensitivity of language models:} The
claim that LLM behaviour is sensitive to contextual
manipulation, making context poisoning a real security
concern. This is supported by documented evidence of
prompt injection and related attacks.

These claims survive the full audit. They exhibit no
significant failure geometries; they are operationally
accessible; their falsification conditions are specifiable;
and their derivation from the framework's mechanisms
is non-trivial and correct.

\begin{auditprop}
\label{ap:level1-stable}
The \LevelI{} claims of the Phoenix corpus constitute
the survivable core $\ThetaS^{(I)}$ of the decomposition
at the engineering level. They are grounded, falsifiable,
non-trivially derived from the framework's mechanisms,
and free of the four failure geometries at significant
severity levels. Their status as survivors is independent
of the audit results for higher-level claims; the
engineering results stand regardless of whether the
philosophical claims can be salvaged.
\end{auditprop}

\section{Level~II Claims: Productive Metaphors}
\label{sec:level2-claims}

\LevelII{} claims are more complex. They are not directly
testable as engineering claims; their value is heuristic
and organizational. But they are not without epistemic
status: a well-chosen analogy generates productive
research directions, reveals structural relationships
between domains, and provides explanatory resources
that pure engineering results do not.

The Phoenix corpus's most productive \LevelII{} claims include:

\textit{Salience as inverse accessibility entropy:}
The structural analogy between Marine jitter metrics
and RSVP accessibility entropy --- low jitter as a
marker of constraint-dense, low-accessibility signal
regions --- is a productive metaphor. It connects the
engineering framework to dynamical systems theory in
a way that suggests specific research directions: testing
whether signals classified as high-salience by the Marine
metric also exhibit lower Kolmogorov complexity, or
occupy more constrained regions of some state space.

\textit{Harmonic consensus as neural population coding:}
The analogy between \HCFone{} conformity encoding and
the efficient coding of correlated neural populations
is productive. It motivates a connection to the
neuroscience of dimensionality reduction and suggests
that the \HCFone{} formalism might provide a computationally
tractable model for testing efficient coding hypotheses.

\textit{Memory as interference:} The \MEMeight{}
architecture as a computational metaphor for the
reconstructive character of biological memory is
productive. It motivates a research connection to
the holographic memory literature and to oscillatory
neuroscience.

\textit{Behavioral injection as semantic malware:}
The analogy between context poisoning in LLMs and
malware operating on reasoning substrates is a
productive metaphor for the AI security literature.
It suggests a research direction in which security
analysis is conducted at the semantic rather than
the executable level.

\begin{auditprop}
\label{ap:level2-productive}
The \LevelII{} claims of the Phoenix corpus constitute
the survivable core $\ThetaS^{(II)}$ at the metaphorical
level. Their value is conditional: they are productive
if treated as analogies that generate research directions
rather than as identity claims that assert the frameworks'
mechanisms are identical across domains. When treated
as identity claims, they migrate to \LevelIV{} and
become subject to the higher standard of ontological
argument.
\end{auditprop}

\section{Level~III Claims: Phenomenological Observations}
\label{sec:level3-claims}

\LevelIII{} claims are phenomenological: they describe
the character of experience rather than the structure
of mechanisms or the performance of implementations.
They are not testable by the standards of engineering
claims, but neither are they mere metaphors; they are
reports on the structure of experience that can be
evaluated against other phenomenological accounts and
against experimental evidence from cognitive science
and psychophysics.

The Phoenix corpus's \LevelIII{} claims include:

\textit{Temporal experience of music:} The claim that
musical experience is constituted by temporal unfolding
rather than spectral content --- that what makes music
meaningful is the pattern of expectation and fulfillment
in time rather than the instantaneous frequency
distribution. This is well-supported in the philosophy
of music and is consistent with psychoacoustic research
on musical tension, expectation, and emotional response.

\textit{Compression degradation as impoverishment:}
The claim that many listeners experience high-quality
compressed audio as impoverished relative to uncompressed
audio, even when they cannot identify the specific
differences. This is a phenomenological report supported
by informal evidence from practitioners and consistent
with (though not established by) the corpus's engineering
proposals.

\textit{Memory as reconstruction:} The claim that
the experience of remembering feels more like reconstruction
than retrieval --- that memories are not experienced
as fixed files accessed from storage but as reconstructed
events that are sensitive to present context. This is
well-established in autobiographical memory research.

These claims are phenomenologically sound and scientifically
supported. Their survival is independent of whether
the mechanisms proposed by the corpus account for them:
the phenomenological observations stand on their own
evidential base.

\begin{auditprop}
\label{ap:level3-phenomenological}
The \LevelIII{} claims of the Phoenix corpus constitute
a partially independent $\ThetaS^{(III)}$: they survive
the decomposition as phenomenological observations that
are independently supported, but their connection to
the corpus's mechanisms requires bridging work that
the corpus has not provided. The existence of the
phenomenological patterns does not establish that the
Marine/\MEMeight{}/\HCFone{} mechanisms account for them.
\end{auditprop}

\section{Level~IV Claims: The Unstable Layer}
\label{sec:level4-claims}

\LevelIV{} claims are ontological: they assert what
cognitive phenomena, consciousness, and meaningful structure
fundamentally are. The Phoenix corpus's \LevelIV{} claims
include:

\textit{Consciousness as temporal coherence:} The claim
that consciousness is constituted by temporal coherence
at the neural level, and that the Marine metric operationalizes
something about consciousness.

\textit{Meaning as temporal structure:} The claim that
meaningful structure \textit{is} temporally coherent
structure --- that there is no meaningful structure
that is not temporally organized in the way the Marine
detects.

\textit{Authenticity as temporal fidelity:} The claim
that authenticity is constituted by temporal structure
preservation, not merely that temporal structure is
one component of authenticity.

These claims are in $\ThetaU$ after the decomposition.
They exhibit high $\AuditP$ scores (the claims about the
underlying space --- consciousness, meaning, authenticity
--- are made on the basis of projected representations
without establishing constraint-faithfulness), high
$\AuditR$ scores (the claims can accommodate any
observation by appeal to coherence dynamics), and
are subject to $\FailGeo{2}$ (the definitional embedding
of the conclusions) and $\FailGeo{3}$ (the broadening
of coherence to universal accommodation).

\begin{auditprop}
\label{ap:level4-unstable}
The \LevelIV{} claims of the Phoenix corpus constitute
the unstable component $\ThetaU^{(IV)}$ after decomposition.
They are not false; they are unestablished in a
specific and diagnosable way: they assert ontological
conclusions at a level that the framework's mechanisms
cannot reach without bridging work that has not been
done and that the failure geometries suggest may not
be achievable within the current framework's structure.
\end{auditprop}

\section{The Level Conflation Problem}
\label{sec:level-conflation}

The most epistemically consequential feature of the
Phoenix corpus's claim structure is not the presence
of claims at \LevelIV{} --- many speculative frameworks
make ontological claims, and some of them turn out to
be correct --- but the conflation of levels within
individual claims and across the corpus's argumentative
structure.

Level conflation occurs when a claim derives its apparent
force from evidence or reasoning appropriate to one
level while making assertions appropriate to a different
level. Several specific instances have been identified:

\textit{The salience-meaning conflation:} Claims that
present evidence from engineering (the Marine detects
X) as supporting ontological conclusions (X is meaningful
structure) without the bridging argument that
would make the inference valid.

\textit{The metaphor-identity conflation:} Arguments
that present structural analogies between the framework's
mechanisms and biological phenomena (\LevelII{}) as
though they established that the mechanisms are the
biological phenomena (\LevelIV{}).

\textit{The phenomenology-mechanism conflation:} Arguments
that present phenomenological observations (\LevelIII{})
as support for specific mechanistic claims (the \MEMeight{}
architecture), when the phenomenological observations
are consistent with many different mechanisms.

Level conflation is the primary mechanism by which the
four failure geometries produce their effects. $\FailGeo{1}$
(Decorative Formalism) makes \LevelI{} notation appear
to be doing \LevelI{} work when it is doing only
organizational work. $\FailGeo{2}$ (Definitional Closure)
makes \LevelII{} definitional choices appear to be
establishing \LevelIV{} ontological conclusions.
$\FailGeo{3}$ (Semantic Universality Collapse) expands
\LevelII{} metaphors until they cover the same ground
as \LevelIV{} claims without the argument. $\FailGeo{4}$
(Operational Severance) names \LevelIV{} quantities
in \LevelI{} technical vocabulary without providing
the measurement procedures that would make the claims
genuinely \LevelI{}.

\medskip

Chapter~14 conducts the RSVP audit: examining the
Relativistic Scalar-Vector Plenum framework against
the same four-level taxonomy and failure geometry
criteria that have been applied to the Phoenix corpus.
This is the most self-critical chapter in Part~III,
because RSVP is the reinterpretation framework that
Part~IV will deploy. Its audit here determines the
conditions under which it can be used --- and the
failure modes it must guard against.
